\newpage
\spread{What the ILM actually reads and writes}\label{page:raw-ilm}
\noindent A transcription can hide omitted strokes, broken words and an incorrect stop. These examples show the image given to the network and the complete image it returned, including blank cells at the end of each grid.
\platehere{Figure_193_raw_ilm}{Original input and output grids from visual inference}{\textbf{a}, The varied-source writer from p.~42 translates the reversal of the request on p.~1. \textbf{b}, The first reserved PAWS-X family in fixed order, shown in both translation directions for the assignment-trained ILM. Both replies stop and lose the statement about a continuous, second-order phase transition. \textbf{c}, Six separate 32-by-32 predictions from the earlier physical-pixel writer, in original development-book order. Each receives the true preceding visual history. All rasters retain their full pixel arrays. Labels lie outside the images.}
\begin{multicols}{2}
\section*{Writing a sequence of glyphs}
The translation ILM reads binary cells in reading order. Each $2\times2$ block contributes its mean ink to a $16\times16$ input. An optical network embeds the cells and the decoder scores its learned visual repertoire and stop at each step. The selected physical glyph supplies the next input to the decoder. The PNG writer lays out the emitted bitmaps, 32 cells to a row. The model can choose an incorrect glyph or an incorrect proposition even when its drawing is sharp.
\section*{Writing pixels within one glyph}
The separate writer in panel c predicts a joint distribution over a cell's binary pixels. Its output can differ from every pictured target in the example. The first six nonempty target positions were chosen before examining their predictions. In the larger qualification, this system marked 25 of 252 true spaces. These conditional cells do not constitute a freely generated sentence. Source and emitted arrays here come directly from the retained evaluation.
\end{multicols}
\noindent{\small \href{https://aclanthology.org/D19-1382/}{PAWS-X, Yang et al. (2019)}. \href{https://github.com/lachlanchen/ImagizedLanguageModel/tree/main/publication/ilm-nature/figures/raw_ilm}{Original PNGs}.}
\newpage
\spread{How much information comes from glyph order?}
\noindent Two passages can contain the same names and objects and describe different events. We trained the writer to detect exchanges of whole translations, then removed both reference images and let it write freely.
\platehere{Figure_194_paired_ink}{A correct pairing can coexist with two incorrect replies}{\textbf{a}, Forty-eight new non-equivalent PAWS-X families, each translated in both directions. Joint pairing compares the two possible assignments. Both row choices requires each source independently to prefer its own complete target. \textbf{b}, Full replies accepted in two opposite, anonymized grading orders. The untouched starting writer also earned zero. Qwen3-8B Q4 received text. The Qwen3-VL evaluator accepted 189 of 192 references and accepted 18 contrasting foils in both orders.}
\begin{multicols}{2}
\section*{The odds of a swap}
Let $s_{ij}=\log p_\theta(Y_j\mid X_i)$ sum over all target images and stop. With equal prior odds for the two assignments and independent conditional draws, the matched-minus-crossed log score is $D=s_{00}+s_{11}-s_{01}-s_{10}$. The normalized probability of a match is $1/(1+e^{-D})$. This is the two-state Gibbs rule with energies equal to negative log likelihoods. Its observed log loss is $\log_2(1+e^{-D})$ bits. A confident wrong match costs more than the one bit charged for a tie.
\par Any additive target fluency or source offset cancels from $D$. If two inputs map to the same representation, $D=0$ for every pair of targets. A deterministic writer must then emit the same image for both. Conversely, arbitrarily large positive $D$ allows a third, incorrect answer to outrank both targets. Lean checks these finite score identities and the collision implication.
\columnbreak
\section*{Reading and writing under intervention}
Both writers saw the same 2,879 training families, ordinary parallel replay and 92,128 directed positive presentations. The added binary loss dropped to 0.0026 nats on training families in the last pass. Held binary loss was 1.648 bits, against 1.651 with likelihood alone and 0.926 for the starting writer. The family-bootstrap interval for the accuracy difference was $[-3.1,3.1]$ percentage points.
\par Sorting each source's unchanged glyphs by bitmap reduced correct pairings from 75/96 to 51/96. Eleven source pairs became identical after sorting. In reading order the model selected the right pairing on all eleven. Sorting produced tied scores and identical free answers on all eleven. Sorting can also disrupt familiar typography and syntax.
\par An inference intervention increased the weight of next-glyph scores that depended on source order. This changed 192/192 PAWS-X and 92/96 OPUS outputs. Neither decoder earned a whole-meaning credit on either panel. The model distinguishes aspects of sequence order, but these interventions have yet to produce faithful independent prose.
\end{multicols}
\noindent{\small Related contrastive evaluations in \href{https://aclanthology.org/2021.blackboxnlp-1.5/}{Vamvas and Sennrich (2021)} and context-aware generation in \href{https://aclanthology.org/2024.naacl-short.69/}{Shi et al. (2024)}.}
